% 封面：南京信息工程大学课程设计报告
% 作为 main.tex 的片段使用：% 封面：南京信息工程大学课程设计报告
% 作为 main.tex 的片段使用：% 封面：南京信息工程大学课程设计报告
% 作为 main.tex 的片段使用：% 封面：南京信息工程大学课程设计报告
% 作为 main.tex 的片段使用：\input{cover}
% 也可用 cover_standalone.tex 单独编译预览
% 所有字段都可在下面几行改：
\def\reporttype{课程设计报告}
\def\reporttitle{基于空域-频域联合分析的深度伪造检测与可解释系统}
\def\studentname{付学瑞}
\def\studentid{202383290286}
\def\major{信息安全}
\def\college{网络空间安全学院}
\def\advisor{}
\def\coversdate{二○二六年十二月}

\begin{titlepage}
\centering
\vspace*{0.9cm}

\includegraphics[width=10.4cm]{uni_name.png}\\[1.0cm]

{\sffamily\bfseries\fontsize{24}{30}\selectfont \reporttype}\\[1.0cm]

\includegraphics[width=3.3cm]{logo.png}\\[1.4cm]

{\fontsize{14}{22}\selectfont
\renewcommand{\arraystretch}{1.75}
\begin{tabular}{r@{\hspace{8pt}}c}
\makebox[3.0cm][r]{\sffamily 题\quad 目} &
  \underline{\makebox[9.1cm][c]{\fontsize{11}{14}\selectfont \reporttitle}} \\[2pt]
\makebox[3.0cm][r]{\sffamily 学生姓名：} &
  \underline{\makebox[9.1cm][c]{\studentname}} \\[2pt]
\makebox[3.0cm][r]{\sffamily 学\qquad 号：} &
  \underline{\makebox[9.1cm][c]{\studentid}} \\[2pt]
\makebox[3.0cm][r]{\sffamily 专\qquad 业：} &
  \underline{\makebox[9.1cm][c]{\major}} \\[2pt]
\makebox[3.0cm][r]{\sffamily 学\qquad 院：} &
  \underline{\makebox[9.1cm][c]{\college}} \\[2pt]
\makebox[3.0cm][r]{\sffamily 指导教师：} &
  \underline{\makebox[9.1cm][c]{\advisor}} \\
\end{tabular}}

\vfill
{\fontsize{15}{20}\selectfont \coversdate}
\vspace{1.0cm}
\end{titlepage}

% 也可用 cover_standalone.tex 单独编译预览
% 所有字段都可在下面几行改：
\def\reporttype{课程设计报告}
\def\reporttitle{基于空域-频域联合分析的深度伪造检测与可解释系统}
\def\studentname{付学瑞}
\def\studentid{202383290286}
\def\major{信息安全}
\def\college{网络空间安全学院}
\def\advisor{}
\def\coversdate{二○二六年十二月}

\begin{titlepage}
\centering
\vspace*{0.9cm}

\includegraphics[width=10.4cm]{uni_name.png}\\[1.0cm]

{\sffamily\bfseries\fontsize{24}{30}\selectfont \reporttype}\\[1.0cm]

\includegraphics[width=3.3cm]{logo.png}\\[1.4cm]

{\fontsize{14}{22}\selectfont
\renewcommand{\arraystretch}{1.75}
\begin{tabular}{r@{\hspace{8pt}}c}
\makebox[3.0cm][r]{\sffamily 题\quad 目} &
  \underline{\makebox[9.1cm][c]{\fontsize{11}{14}\selectfont \reporttitle}} \\[2pt]
\makebox[3.0cm][r]{\sffamily 学生姓名：} &
  \underline{\makebox[9.1cm][c]{\studentname}} \\[2pt]
\makebox[3.0cm][r]{\sffamily 学\qquad 号：} &
  \underline{\makebox[9.1cm][c]{\studentid}} \\[2pt]
\makebox[3.0cm][r]{\sffamily 专\qquad 业：} &
  \underline{\makebox[9.1cm][c]{\major}} \\[2pt]
\makebox[3.0cm][r]{\sffamily 学\qquad 院：} &
  \underline{\makebox[9.1cm][c]{\college}} \\[2pt]
\makebox[3.0cm][r]{\sffamily 指导教师：} &
  \underline{\makebox[9.1cm][c]{\advisor}} \\
\end{tabular}}

\vfill
{\fontsize{15}{20}\selectfont \coversdate}
\vspace{1.0cm}
\end{titlepage}

% 也可用 cover_standalone.tex 单独编译预览
% 所有字段都可在下面几行改：
\def\reporttype{课程设计报告}
\def\reporttitle{基于空域-频域联合分析的深度伪造检测与可解释系统}
\def\studentname{付学瑞}
\def\studentid{202383290286}
\def\major{信息安全}
\def\college{网络空间安全学院}
\def\advisor{}
\def\coversdate{二○二六年十二月}

\begin{titlepage}
\centering
\vspace*{0.9cm}

\includegraphics[width=10.4cm]{uni_name.png}\\[1.0cm]

{\sffamily\bfseries\fontsize{24}{30}\selectfont \reporttype}\\[1.0cm]

\includegraphics[width=3.3cm]{logo.png}\\[1.4cm]

{\fontsize{14}{22}\selectfont
\renewcommand{\arraystretch}{1.75}
\begin{tabular}{r@{\hspace{8pt}}c}
\makebox[3.0cm][r]{\sffamily 题\quad 目} &
  \underline{\makebox[9.1cm][c]{\fontsize{11}{14}\selectfont \reporttitle}} \\[2pt]
\makebox[3.0cm][r]{\sffamily 学生姓名：} &
  \underline{\makebox[9.1cm][c]{\studentname}} \\[2pt]
\makebox[3.0cm][r]{\sffamily 学\qquad 号：} &
  \underline{\makebox[9.1cm][c]{\studentid}} \\[2pt]
\makebox[3.0cm][r]{\sffamily 专\qquad 业：} &
  \underline{\makebox[9.1cm][c]{\major}} \\[2pt]
\makebox[3.0cm][r]{\sffamily 学\qquad 院：} &
  \underline{\makebox[9.1cm][c]{\college}} \\[2pt]
\makebox[3.0cm][r]{\sffamily 指导教师：} &
  \underline{\makebox[9.1cm][c]{\advisor}} \\
\end{tabular}}

\vfill
{\fontsize{15}{20}\selectfont \coversdate}
\vspace{1.0cm}
\end{titlepage}

% 也可用 cover_standalone.tex 单独编译预览
% 所有字段都可在下面几行改：
\def\reporttype{课程设计报告}
\def\reporttitle{基于空域-频域联合分析的深度伪造检测与可解释系统}
\def\studentname{付学瑞}
\def\studentid{202383290286}
\def\major{信息安全}
\def\college{网络空间安全学院}
\def\advisor{}
\def\coversdate{二○二六年十二月}

\begin{titlepage}
\centering
\vspace*{0.9cm}

\includegraphics[width=10.4cm]{uni_name.png}\\[1.0cm]

{\sffamily\bfseries\fontsize{24}{30}\selectfont \reporttype}\\[1.0cm]

\includegraphics[width=3.3cm]{logo.png}\\[1.4cm]

{\fontsize{14}{22}\selectfont
\renewcommand{\arraystretch}{1.75}
\begin{tabular}{r@{\hspace{8pt}}c}
\makebox[3.0cm][r]{\sffamily 题\quad 目} &
  \underline{\makebox[9.1cm][c]{\fontsize{11}{14}\selectfont \reporttitle}} \\[2pt]
\makebox[3.0cm][r]{\sffamily 学生姓名：} &
  \underline{\makebox[9.1cm][c]{\studentname}} \\[2pt]
\makebox[3.0cm][r]{\sffamily 学\qquad 号：} &
  \underline{\makebox[9.1cm][c]{\studentid}} \\[2pt]
\makebox[3.0cm][r]{\sffamily 专\qquad 业：} &
  \underline{\makebox[9.1cm][c]{\major}} \\[2pt]
\makebox[3.0cm][r]{\sffamily 学\qquad 院：} &
  \underline{\makebox[9.1cm][c]{\college}} \\[2pt]
\makebox[3.0cm][r]{\sffamily 指导教师：} &
  \underline{\makebox[9.1cm][c]{\advisor}} \\
\end{tabular}}

\vfill
{\fontsize{15}{20}\selectfont \coversdate}
\vspace{1.0cm}
\end{titlepage}
